\section{Introduction}

A negative result has at least three possible meanings: the candidate mechanism is wrong, the current representation/search family is inadequate, or the computation could not establish either conclusion. Research software commonly collapses these states. A failed run may be treated as negative evidence; a new idea may silently replace the comparison set; a donor-equivalent method may be described as new; or an old negative may be reopened without recording which assumption changed.

This paper studies a narrower object than a general theory of scientific discovery: an executable discipline for preserving and recursively revisiting negative computational results. The discipline is built around five constraints. Outcome rules are frozen before result access. Donor methods receive first right of refusal. Exact claims carry witnesses or independently checkable referees. Failure and authority are typed so a capability result cannot self-authorize a scientific conclusion. Finally, the record is replayable enough that a later reader can reconstruct what was known, what was attempted, and why a lane stopped.

The programme studied here provides a stress test because its history contains the outcomes that a non-gamed method should retain: donor absorptions, exact refutations, bounded positives, mixed results, and unresolved/cannot-check states. Its strongest mathematical result, its dual-instrument decision benchmark, and its partial-knowledge mechanism studies are each reported separately and are cited here as ordinary references. They matter to this paper only as evidence that the common methodology did not force them into one preferred terminal.

The central claim is therefore descriptive and methodological: this receipt-first negative-recovery workflow was instantiated end-to-end and behaved in a non-monotone, outcome-sensitive way under prospectively frozen rules. We do not estimate how often such a workflow improves science, and we do not claim that receipt completeness is a substitute for evidentiary quality.